\chapter{Acupuncture}

\section*{Overview}
Acupuncture is a component of traditional Chinese medicine in which thin sterile needles are inserted into specific points on the body to modulate symptoms. It is most commonly used in the West for pain relief and for nausea and vomiting.

\section*{Alternative Names}
Traditional acupuncture; needling therapy; acupressure (needleless variant); auricular acupuncture

\section*{Principle}
According to traditional theory, acupuncture restores the flow of vital energy (qi) along meridians to rebalance the body. From a biomedical perspective, needle insertion at acupuncture points stimulates local and central nervous system pathways, releasing endogenous opioids, serotonin, and other neuromodulators that reduce pain and influence autonomic function.

\section*{History}
Acupuncture originated in ancient China more than two millennia ago and was systematized during the Han dynasty, spreading to Europe and the Americas from the 17th century onward.

\section*{Why This Test or Procedure Is Ordered}
Ordered or recommended for chronic pain conditions such as low back pain, osteoarthritis, and neck pain, and for postoperative and chemotherapy-related nausea and vomiting. It is also used for migraine, fibromyalgia, and as an adjunct in other conditions in which patients seek nonpharmacologic therapy.

\section*{Is This a Primary or Secondary Test or Procedure}
A secondary or complementary therapy used alongside, or when conventional treatment is insufficient.

\section*{How This Test or Procedure Is Performed}
After a history and examination, the practitioner inserts sterile single-use needles at selected points to a depth of a few millimeters to a few centimeters and may briefly twist or electrically stimulate them. The needles are left in place for roughly 10 to 30 minutes while the patient lies still, and the session ends with removal and disposal of the needles.

\section*{What Preparation Is Needed}
No special preparation is needed. The patient should eat lightly and wear loose clothing, and should inform the practitioner of bleeding disorders, anticoagulant use, and pregnancy.

\section*{Contraindications and Precautions}
Acupuncture is generally avoided at sites of skin infection, scar tissue, or lymphedema, and needle use is avoided over a pacemaker site or during electrical stimulation in patients with implanted devices. Patients with severe bleeding disorders or on anticoagulation should be treated cautiously, and those with needle phobia may not tolerate the procedure.

\section*{Risks}
Serious complications are rare when performed by a qualified practitioner. Mild bruising, soreness, or bleeding at the site may occur, and rare complications include infection, pneumothorax, and organ puncture.

\section*{What Happens After the Test or Procedure}
The patient may feel relaxed or mildly sore and can resume normal activities immediately. Repeat sessions are scheduled according to the treatment plan, with symptom response reviewed periodically.

\section*{Understanding Your Results}
Improvement is judged by symptom scores, pain reduction, and functional gains over a course of several sessions. If no benefit is seen after a reasonable number of treatments, further therapy is generally not recommended.

\section*{Normal Results}
Relief or meaningful reduction in the presenting symptom without significant adverse effects.

\section*{Follow-up Tests and Procedures}
Symptom diaries and pain scales are used to track response between visits. If the response is incomplete or symptoms persist, the treatment plan is revised or additional evaluation for an underlying cause is pursued.

\section*{References}
\begin{enumerate}
  \item MedlinePlus. Acupuncture. U.S. National Library of Medicine; 2025.
  \item Current Medical Diagnosis and Treatment. McGraw-Hill; 2025.
\end{enumerate}

\clearpage
